% =========================================================================
%  Figura (Problema 3.5): placas planas paralelas.
%  (a) Condensador ideal: +sigma y -sigma en las caras enfrentadas, campo
%      sigma/eps0 entre las placas y nulo fuera; caja de Gauss.
%  (b) Cargas cualesquiera: cuatro caras q1..q4 y campo en cada región.
% =========================================================================
\begin{figure}[t]
\centering
\begin{tikzpicture}[font=\small]
  % ---------------------------------------------------------------- (a)
  \begin{scope}
    \fill[top color=azulvelo!80!white, bottom color=metal, draw=azulprof, line width=0.6pt] (-1.75,-1.7) rectangle (-1.4,1.7);
    \fill[top color=azulvelo!80!white, bottom color=metal, draw=azulprof, line width=0.6pt] (1.4,-1.7) rectangle (1.75,1.7);
    \foreach \y in {-1.5,-1.1,...,1.6}{\node[signopos, font=\tiny\bfseries] at (-1.28,\y) {$+$};
                                       \node[signoneg, font=\tiny\bfseries] at (1.28,\y) {$-$};}
    \foreach \y in {-1.35,-0.85,0.85,1.35}{\draw[lcampo] (-1.1,\y) -- (1.1,\y);}
    \node[font=\footnotesize, fill=white, inner sep=1.5pt] at (0.25,-0.2) {$E=\dfrac{\sigma}{\eo}$};
    \draw[gaussazul] (-1.62,0.15) rectangle (-0.75,0.55);
    \node[font=\scriptsize, text=azulacero, anchor=west] at (-0.7,0.35) {gaussiana};
    \node[font=\footnotesize] at (-2.35,-0.9) {$E=0$};
    \node[font=\footnotesize] at (2.35,-0.9) {$E=0$};
    \node[font=\footnotesize, text=cargapos] at (-1.58,-2.0) {$+Q$};
    \node[font=\footnotesize, text=carganeg] at (1.58,-2.0) {$-Q$};
    \draw[{Stealth[length=3pt]}-{Stealth[length=3pt]}, grispizarra] (-1.4,1.95) -- (1.4,1.95) node[midway, above, font=\footnotesize] {$d$};
    \node at (0,-3.15) {(a)};
  \end{scope}
  % ---------------------------------------------------------------- (b)
  \begin{scope}[xshift=8.6cm]
    \fill[top color=azulvelo!80!white, bottom color=metal, draw=azulprof, line width=0.6pt] (-2.2,-1.7) rectangle (-1.6,1.7);
    \fill[top color=azulvelo!80!white, bottom color=metal, draw=azulprof, line width=0.6pt] (1.6,-1.7) rectangle (2.2,1.7);
    \foreach \x/\q in {-2.2/q_1,-1.6/q_2,1.6/q_3,2.2/q_4}{
      \draw[cargapos, line width=1.3pt] (\x,-1.7) -- (\x,1.7);
      \node[font=\footnotesize, anchor=south] at (\x,1.72) {$\q$};}
    \node[font=\scriptsize, rotate=90] at (-1.9,0) {$\vE=\vc{0}$};
    \node[font=\scriptsize, rotate=90] at (1.9,0) {$\vE=\vc{0}$};
    \node[font=\footnotesize] at (-1.9,-2.0) {$Q_1$};
    \node[font=\footnotesize] at (1.9,-2.0) {$Q_2$};
    \draw[-{Stealth[length=4pt]}, azulacero, line width=0.8pt] (-2.45,1.0) -- (-3.25,1.0);
    \node[font=\scriptsize, anchor=north] at (-2.85,0.85) {$\dfrac{Q_1+Q_2}{2\eo A}$};
    \draw[-{Stealth[length=4pt]}, azulacero, line width=0.8pt] (-0.4,1.0) -- (0.4,1.0);
    \node[font=\scriptsize, anchor=north] at (0,0.85) {$\dfrac{Q_1-Q_2}{2\eo A}$};
    \draw[-{Stealth[length=4pt]}, azulacero, line width=0.8pt] (2.45,1.0) -- (3.25,1.0);
    \node[font=\scriptsize, anchor=north] at (2.85,0.85) {$\dfrac{Q_1+Q_2}{2\eo A}$};
    \node[font=\scriptsize, align=center] at (0,-2.45) {$q_1=q_4=\tfrac12(Q_1+Q_2)$\\[3pt] $q_2=-q_3=\tfrac12(Q_1-Q_2)$};
    \node at (0,-3.15) {(b)};
  \end{scope}
\end{tikzpicture}
\caption{Problema 3.5. (a) Condensador plano ideal: la carga está en las caras enfrentadas y el campo, uniforme, solo existe entre las placas; la caja de Gauss
con una tapa dentro del metal da $E=\sigma/\eo$. (b) Placas con cargas cualesquiera: las caras exteriores se reparten la carga total y las interiores llevan
cargas opuestas; se indica el campo en cada región (para $Q_1>Q_2$ y $Q_1+Q_2>0$).}
\label{fig:sem-plano}
\end{figure}
